\documentclass[12pt]{article}
\newcommand{\bottomMargin}{2cm}


\newcommand{\header}{
	ДИЗАЙН И АНАЛИЗ НА АЛГОРИТМИ - \\ ПРАКТИКУМ\\
	\vspace{0.1cm}
	Летен семестър, 2026 г., първо контролно\\
	\vspace{0.1cm}
}
\newcommand{\tl}{$0.5$ сек.}
\newcommand{\ml}{$256$ MB}

\input{structure.tex}

\begin{document}
	
	\section{Задача К4. Боксиране}
	
	
	Ще се организира мач по бокс по повод утрешния рожден ден на Сашка. За целта ще бъдат избрани двама боксьори. Има точно $N(N+1)$ боксьори, които искат да участват.
	
	Боксьорите тренират на трасе от $N$ боксови круши, номерирани с числа от $1$ до $N$ спрямо позицията им в трасето отляво-надясно, като $i$-тата от тях има сила $a_i$. За всяка двойка числа $(l,r)$, за които $1 \leq l \leq r \leq N$, има двойка боксьори, които държат да тренират на подтрасето, образувано от всички боксови круши с номера от $l$ до $r$, като се запази подредбата им от първоначалното трасе. 
	
	Начинът, по който тренират двойка боксьори, е много интересен. Нека дадена двойка да тренира върху подтрасе с круши, номерирани с числата от 1 до $M$ спрямо позицията им. Двамата боксьори ще изберат някаква позиция $idx$ $(0 \leq idx \leq M)$, като единият от тях ще тренира с всяка круша от $1$ до $idx$, а другият – с всяка круша от $idx+1$ до $M$. Възможно е някой от двамата да не тренира с нито една боксова круша. Всяка круша в подтрасето трябва да бъде използвана от точно един от боксьорите. Всеки от боксьорите ще получи подготовка, равна на сбора от силите на боксовите круши, с които е тренирал. Ако един от двамата боксьори не е тренирал с нито една круша, неговата подготвеност ще е равна на $0$. 
	
	На Сашка ще $\grave{\text{и}}$ направят впечатление тези двойки боксьори, в които и двамата играчи имат подготовка, по-голяма или равна на $K$. Организаторите пък се притесняват за състоянието на боксовите круши, но въпреки това искат двойката боксьори, които ще изберат, да направи впечатление на Сашка. Заради това Ви молят да напишете програма \texttt{\textbf{boxing}}, която да намери минималния брой последователни круши, отговарящи на трасе на двойка боксьори, които ще впечатлят Сашка. Ако не съществува трасе на двойка боксьори, които ще впечатлят Сашка, изведете $-1$.
	
	\subsection{Вход}
	Тъй като броят на крушите $N$ може да е много голям, техните сили няма да бъдат подадени директно на входа. Вместо това, вие трябва да ги генерирате сами в програмата си. 
	
	На първия ред ще получите числата $N$ и $K$. На втория ред ще бъдат дадени $a_1, a_2, X, Y, Z$ и $M$.
	
	\begin{itemize}
		\item Първите два елемента на редицата са съответно $a_1$ и $a_2$.
		\item Всеки следващ елемент (за $i \geq 3$) се пресмята по формулата: 
		$a_i = (a_{i-2} \cdot X + a_{i-1} \cdot Y + Z) \bmod M$
	\end{itemize}
	
	\textit{(Забележка: Ако $N=2$, просто игнорирайте променливите за генериране).}
	
	\subsection{Изход}
	На стандартния изход отпечатайте търсения минимален брой последователни круши.
	
	\subsection{Ограничения}
	\vspace{0.1em}
	\begin{itemize}
		\item $1 \leq N \leq 10^7$
		\item $1 \leq K \leq 10^{16}$
		\item $0 \leq a_1,a_2 < M$
		\item $1 \leq X,Y,Z \leq M \leq 2.10^9$
	\end{itemize}
	
	\subsection{Оценяване}
	\noindent В $30\%$ от тестовете $N \leq 10 000$. \\
	В други тестове, носещи около $20\%$ от точките: $N \leq 10^5$.\\
	В други тестове, носещи около $20\%$ от точките: $N \leq 2.10^6$.\\
	В останалите тестове, носещи около $30\%$ от точките, няма допълнителни
	ограничения.
	
	\subsection{Примери}
	\begin{table}[ht]
		\begin{tblr}{|X[42,l]|X[8,l]|X[50,l]|}
			\hline
			\textbf{Вход} & \textbf{Изход} & \textbf{Обяснение на примера} \\
			\hline
			\texttt{2 11\\10 6 12 15 13 15} & 
			\texttt{-1} &
			{Редицата е \texttt{10 6}.}\\
			\hline
			\texttt{5 13\\10 6 12 15 13 15} & 
			\texttt{3} &
			{Редицата е \texttt{10 6 13 10 4}} \\
			\hline
			\texttt{50 6012\\612 419 836 700 716 1009} & 
			\texttt{20} &
			{Редицата е \texttt{612 419 466 161 510 929 773 706 \dots}} \\
			\hline
		\end{tblr}
	\end{table}
	
\end{document}
